% !PRINTED 567
\weberpage{567}
Da die $f$-gliedrigen Perioden (f\"ur $f = 1$) die Gr\"ossen $r$ mit
umfassen, so wollen wir gleich zur Bestimmung dieser Perioden
\"ubergehen.

Wir wollen also unter $\alpha$ eine primitive $e^{\mathrm{te}}$ Einheitswurzel
verstehen und $\mu = f$ setzen.
Dann ist
\[ (\alpha^\lambda, r) = r + \alpha^\lambda r_1 + \alpha^{2\lambda} r_2 + \cdots + \alpha^{(n-2)\lambda} r_{n-2}, \]
oder nach \S.\,167, (12)
\weq{11}{(\alpha^\lambda, r) = (\alpha^\lambda, \eta) = \eta + \alpha^\lambda \eta_1 + \alpha^{2\lambda} \eta_2 + \cdots + \alpha^{(e-1)\lambda} \eta_{e-1},}
und die Formeln (9) und (10) ergeben
\weq{12}{(\alpha^\lambda, \eta)(\alpha, \eta) = (\alpha^{\lambda+1}, \eta)\psi_\lambda(\alpha) \qquad \lambda = 1, 2 \ldots (e-2),}
\weq{13}{(\alpha^\lambda, \eta)(\alpha^{-\lambda}, \eta) = (-1)^{f\lambda} n.}
Die Gleichung (13) zeigt, dass keine der Zahlen $(\alpha^\lambda, \eta)$ verschwindet,
und nach (12) sind auch die $\psi_\lambda(\alpha)$ von Null verschieden,
und wenn wir die Gleichung (12) f\"ur $\lambda = 1, 2 \ldots \lambda - 1$
bilden, so folgt
% [source-original mathematical notation, verified 2026-08-29] The print uses lambda
% both for the running values and for the fixed upper index in "lambda-1". The displayed
% sequence below makes the intended range clear; the diplomatic text is preserved.
\[
\begin{aligned}
(\alpha, \eta)(\alpha, \eta) &= (\alpha^2, \eta)\psi_1(\alpha)\\
(\alpha^2, \eta)(\alpha, \eta) &= (\alpha^3, \eta)\psi_2(\alpha)\\
&\cdots\cdots\cdots\cdots\cdots\cdots\cdots\cdots\cdots\cdots\\
(\alpha^{\lambda-1}, \eta)(\alpha, \eta) &= (\alpha^\lambda, \eta)\psi_{\lambda-1}(\alpha).
\end{aligned}
\]
Daraus erhalten wir durch Multiplication und Wegheben
des Factors $(\alpha^2, \eta) \ldots (\alpha^{\lambda-1}, \eta)$
\weq{14}{(\alpha, \eta)^\lambda = (\alpha^\lambda, \eta)\psi_1(\alpha)\psi_2(\alpha) \ldots \psi_{\lambda-1}(\alpha).}
Dadurch ist $(\alpha^\lambda, \eta)$ rational durch $(\alpha, \eta)$ ausgedr\"uckt. Setzen
wir aber in (14) $\lambda = e - 1$, so ergiebt sich
\[ (\alpha, \eta)^{e-1} = (\alpha^{-1}, \eta)\psi_1(\alpha)\psi_2(\alpha) \ldots \psi_{e-2}(\alpha), \]
und durch Multiplication mit $(\alpha, \eta)$, wenn (13) in der Form
\[ (\alpha, \eta)(\alpha^{-1}, \eta) = (-1)^f n \]
benutzt wird,
\weq{15}{(\alpha, \eta)^e = (-1)^f n\, \psi_1(\alpha)\psi_2(\alpha) \ldots \psi_{e-2}(\alpha),}
wodurch die Bildung von $(\alpha, \eta)$ auf eine $e^{\mathrm{te}}$ Wurzel zur\"uckgef\"uhrt ist. Setzen wir $e = m$, so erhalten wir $(\varepsilon, r)$.

Wir wollen dies Verfahren auf den interessanten Fall $n = 17$
anwenden, der zu dem geometrisch merkw\"urdigen, von \gesperrt{Gauss}
zuerst gefundenen Resultat f\"uhrt, dass die Theilung der Kreisperipherie
 in 17 Theile von einer Reihe quadratischer Gleichungen
 abh\"angt, und dass daher das regul\"are Siebzehneck mit
Cirkel und Lineal construirt werden kann.
